\chapter{时间序列与统计套利：从动态模型到滚动验证}

\begin{itemize}
  \item \textbf{直接先修：}上册第 9、11、12、14、15、16、17 章；诊断入口见第 2.3 节。
  \item \textbf{学习产出：}交付一个能复现动态模型、状态估计、变点、成本台账与失败边界的滚动研究包。
  \item \textbf{研究边界：}本章证明的是合成协议和诊断链路，不把历史相关性宣称为稳定套利关系。
\end{itemize}

\section{ARMA、VAR 与 GARCH：先写信息集}

ARMA 模型描述条件均值，GARCH 描述条件方差，VAR 则把多个序列的滞后共同放入均值方程。任何模型都应先声明信息集 $\mathcal F_t$、采样频率、缺失处理与预测时点。例如 AR(1)
\[
X_t=\phi X_{t-1}+\varepsilon_t,
\qquad \mathbb E(\varepsilon_t\mid\mathcal F_{t-1})=0,
\]
只有在弱平稳、创新条件和 $|\phi|<1$ 等假设下，才有稳定的二阶解释。VAR 还需报告滞后阶、稳定根和变量排序；GARCH 需要非负参数、平稳条件和标准化残差诊断。

已知路径 $1,0.6,0.36,0.216,0.1296$ 满足 $X_t=0.6X_{t-1}$，故无截距最小二乘
\[
\widehat\phi=\frac{\sum_{t=2}^T X_{t-1}X_t}{\sum_{t=2}^T X_{t-1}^2}=0.6.
\]
若错误地拟合“均值为零、无动态”的模型，残差就是原序列，其一阶回归系数仍为 $0.6$，超过冻结阈值 $0.2$，必须拒绝。一般残差自相关写作 $\rho(h)=\operatorname{Corr}(e_t,e_{t-h})$。Ljung--Box、残差 ACF 与标准化残差图的作用是暴露错设，而不是为已选模型盖章。

\begin{diagnostic}
随机打乱时间再交叉验证会破坏滞后关系；用未来波动率标准化当前收益则直接改变 $\mathcal F_t$。预测分数再高也不能恢复被破坏的时间协议。
\end{diagnostic}

\section{协整不是高相关}

两个 $I(1)$ 序列 $x_t,y_t$ 若存在 $\beta$ 使
\[
z_t=y_t-\beta x_t\sim I(0),
\]
才称协整。高相关可能只是共同趋势；低短期相关也不排除稳定误差修正关系。本章小样本取 $\beta=2$，手算 $z_t\in\{-1,1\}$ 且均值为零，用它核对实现。完整研究仍需单位根前检验、确定性项、结构突变和样本外稳定性。

误差修正模型把短期变化与长期偏离分开：
\[
\Delta y_t=\alpha z_{t-1}+\sum_j\Gamma_j\Delta w_{t-j}+u_t.
\]
$\alpha$ 的符号和速度只有在协整向量与时点协议冻结后才有意义。先浏览全样本选“最稳定”的配对，再回到起点交易，是伪配对；把未来样本估计的 $\beta$ 用于过去也属于泄漏。

\section{状态空间与手算 Kalman 更新}

线性高斯状态空间写成
\[
s_t=Fs_{t-1}+\eta_t,\qquad y_t=Hs_t+\epsilon_t.
\]
标量更新中，先验均值 $0$、方差 $1$，观测 $1.2$、观测噪声方差 $0.25$，故
\[
K=\frac{1}{1+0.25}=0.8,
\quad m^+=0+0.8(1.2)=0.96,
\quad P^+=(1-0.8)1=0.2.
\]
这三个数是独立 oracle。生产实现还要记录预测步、更新步、缺测规则、初始分布与数值分解。用平滑器得到的未来状态解释过去交易，等同把未来观测送入在线滤波，门禁必须拒绝。

\begin{note}[Python 边界]
Python 库常把滤波、平滑和缺测处理封装在对象里；数学上它们对应不同条件信息。回测只能使用 $p(s_t\mid y_{1:t})$，不能悄悄替换成 $p(s_t\mid y_{1:T})$。
\end{note}

\section{状态切换、变点与可识别性}

状态切换模型需要说明状态标签约束、转移矩阵、发射分布与初值；否则标签交换会让“牛市状态 1”之类的解释失去可识别性。变点检测则必须预先冻结最短区间、阈值和损失函数。

合成序列前四项围绕 0，后四项围绕 3。最大化两段均值差得到变点索引 4；在冻结阈值下检测延迟为 0、提前误报为 0。报告必须同时给出延迟、误报和阈值敏感性，不能只展示命中的那个案例。若把全样本变点反向用于训练前的仓位，仍然是未来信息。

\section{Walk-forward 与交易台账}

本章冻结三个互不重叠的窗口：训练 2016--2019，验证 2020--2021，交易 2022--2024。每个滚动步只能用过去拟合、紧随其后的验证选择、再进入尚未打开的交易段。随机切分和全样本标准化统一返回稳定拒绝类别。

小样本两笔价差毛收益为 $0.03$ 与 $0.02$，总成本 $0.01$，所以
\[
R^{\mathrm{net}}=0.03+0.02-0.01=0.04.
\]
真实台账还应逐笔保存价差、滑点、佣金、借券、持有期、未成交与强平；净收益必须由交易台账独立重算，而不是从策略汇总表抄取。

\begin{implementationnote}
\begin{lstlisting}[language=bash]
uv run python notebooks/lower/ch02_stat_arb.py \
  evidence/lower-ch02/oracle.json
\end{lstlisting}
命令逐项核对 AR 错设、Kalman 更新、变点、协整价差、walk-forward、成本与三种失败门禁。
\end{implementationnote}

\section{Route Capstone：滚动统计套利研究包}

Capstone 需保留原始滚动结果，而不是只保存最终曲线。目录至少含：冻结输入与哈希、每窗参数与残差、状态概率或变点结果、交易台账、成本假设、失败注入和限制报告。独立基线可以是无状态的固定阈值策略；复杂模型必须证明在相同时间边界和成本下增加了什么。

冻结摘要见 \path{reports/lower-ch02-summary.md}。报告应讨论参数漂移、协整失效、检测延迟、误报、成本敏感性与跨市场差异，并明确写出：历史相关性不等于稳定套利关系。

\section{四级问题}
\begin{enumerate}[leftmargin=*]
  \item \textbf{口述概念：}白噪声、创新和错设残差为什么不是同一对象？
  \item \textbf{口述概念：}高相关、协整与可交易价差分别需要哪些额外条件？
  \item \textbf{笔试推导：}由几何 AR 路径推导 $\widehat\phi=0.6$，并解释错设诊断。
  \item \textbf{笔试推导：}手算标量 Kalman 增益、后验均值与方差。
  \item \textbf{数值编程：}改变变点阈值，报告延迟、误报和未检出三类结果。
  \item \textbf{数值编程：}增加第三笔交易，独立重算毛收益、成本和净收益。
  \item \textbf{研究判断：}为什么随机切分和全样本标准化会夸大统计套利表现？
  \item \textbf{研究判断：}设计一个协整失效后的停用协议，并说明它不能回看未来。
\end{enumerate}

